\begingroup
\renewcommand{\clearpageforsection}{}
\exercisesheader{}
\endgroup


% 31

\eoce{\qt{Benar atau salah, Bagian I\label{tf_chisq_1}} Tentukan apakah pernyataan di bawah ini 
benar atau salah. Untuk setiap pernyataan yang salah, usulkan rumusan alternatif agar 
pernyataan tersebut menjadi benar.
\begin{parts}
\item Distribusi khi-kuadrat, sama seperti distribusi normal, memiliki dua 
parameter, yaitu rata-rata dan simpangan baku.
\item Distribusi khi-kuadrat selalu menceng ke kanan, terlepas dari 
nilai parameter derajat kebebasan.
\item Statistik khi-kuadrat selalu positif.
\item Ketika derajat kebebasan meningkat, bentuk distribusi khi-kuadrat 
menjadi semakin menceng.
\end{parts}
}{}

% 32

\eoce{\qt{Benar atau salah, Bagian II\label{tf_chisq_2}} Tentukan apakah pernyataan di bawah ini 
benar atau salah. Untuk setiap pernyataan yang salah, usulkan rumusan alternatif agar 
pernyataan tersebut menjadi benar.
\begin{parts}
\item Ketika derajat kebebasan meningkat, rata-rata distribusi khi-kuadrat 
meningkat.
\item Jika Anda memperoleh $\chi^2 = 10$ dengan $df = 5$, Anda akan gagal menolak $H_0$ 
pada tingkat signifikansi 5\%.
\item Ketika mencari nilai-p dari uji khi-kuadrat, kita selalu mengarsir daerah 
ekor pada kedua ekor.
\item Ketika derajat kebebasan meningkat, keragaman distribusi khi-kuadrat 
menurun.
\end{parts}
}{}

% 33

\eoce{\qt{Buku teks sumber terbuka\label{opensource_text_chisq_GOF}} Seorang dosen yang menggunakan 
buku pengantar statistika sumber terbuka memperkirakan bahwa 60\% 
mahasiswa akan membeli salinan cetak buku tersebut, 25\% akan mencetaknya dari 
web, dan 15\% akan membacanya secara daring. Pada akhir semester, dosen tersebut meminta 
mahasiswanya mengisi survei yang menanyakan format buku 
yang mereka gunakan. Dari 126 mahasiswa, 71 menyatakan membeli salinan cetak buku, 
30 menyatakan mencetaknya dari web, dan 25 menyatakan membacanya secara daring.
\begin{parts}
\item Nyatakan hipotesis untuk menguji apakah prediksi dosen tersebut 
tidak akurat.
\item Berapa banyak mahasiswa yang diperkirakan dosen akan membeli buku, mencetak 
buku, dan membaca buku hanya secara daring?
\item Situasi ini sesuai untuk uji khi-kuadrat. Daftarkan 
syarat-syarat yang diperlukan untuk uji tersebut dan periksa bahwa semuanya terpenuhi.
\item Hitung statistik khi-kuadrat, derajat kebebasan yang terkait 
dengannya, dan nilai-p.
\item Berdasarkan nilai-p yang dihitung pada bagian (d), apa kesimpulan 
uji hipotesis tersebut? Tafsirkan kesimpulan Anda dalam konteks ini.
\end{parts}
}{}

% 34

\eoce{\qt{Kijang\label{barking_deer_chisq_GOF}}
Faktor mikrohabitat yang berkaitan dengan lokasi mencari makan dan tempat beristirahat
kijang di Pulau Hainan, Tiongkok, diteliti.
Di wilayah ini, hutan mencakup 4.8\% dari luas daratan,
petak rumput budidaya mencakup 14.7\%, dan hutan gugur
mencakup 39.6\%.
Dari 426 lokasi tempat kijang mencari makan, 4 dikategorikan
sebagai hutan, 16 sebagai petak rumput budidaya, dan 61 sebagai hutan gugur.
Tabel di bawah ini merangkum data tersebut.\footfullcite{Teng:2004}
\begin{center}
\begin{tabular}{c c c c c}
Hutan	& Rumput budidaya	& Hutan gugur	 & Lainnya & Total \\
\hline 
4		& 16					& 61			     & 345	 & 426 \\
\end{tabular}
\end{center}

\noindent \begin{minipage}[c]{0.7\textwidth}
\begin{parts}
\item Tuliskan hipotesis untuk menguji apakah kijang lebih memilih mencari makan di 
habitat tertentu daripada habitat lainnya.
\item Jenis uji apa yang dapat kita gunakan untuk menjawab pertanyaan penelitian ini?
\item Periksa apakah asumsi dan syarat yang diperlukan untuk uji tersebut 
terpenuhi.
\item Apakah data ini memberikan bukti meyakinkan bahwa kijang lebih memilih 
mencari makan di habitat tertentu daripada habitat lainnya? Lakukan uji hipotesis 
yang sesuai untuk menjawab pertanyaan penelitian ini.
\end{parts}
\end{minipage}
\begin{minipage}[c]{0.03\textwidth}
$\:$ \\
\end{minipage}
\begin{minipage}[c]{0.28\textwidth}
\begin{center}
\Figures[Foto seekor kijang bertanduk bercabang dan berwarna cokelat kemerahan, mengintip dari sela daun dan vegetasi.]{0.7}{eoce/barking_deer_chisq_GOF}{barking_deer.jpg} \\
{\footnotesize Foto oleh Shrikant Rao (\oiRedirect{textbook-flickr_shrikant_rao_barking_deer}{http://flic.kr/p/4Xjdkk}) \oiRedirect{textbook-CC_BY_2}{lisensi CC~BY~2.0}}
\end{center}
\end{minipage}
}{}
